% section: 反復・不動点・非線形漸化式
% 移行元:
%   chapter_04_appendix/section_03_dynamics/section.tex
%   chapter_04_appendix/section_10_mobius_riccati/section.tex
%   chapter_04_appendix/section_08_special_functions_bridge/section.tex の翻訳部分

\section{反復・不動点・非線形漸化式}

第2章では、漸化式を式変形によって既知の型へ帰着した。
ここでは、漸化式そのものを関数の反復として見る。
この見方を使うと、一般項を求めることが難しい場合でも、数列がどこへ進むのか、
どのような点を中心に動くのかを調べられる。

\subsection{数列を関数の反復として見る}

一項前から次の項を決める漸化式
\[
  a_{n+1}=f(a_n)
\]
を考える。初期値 $a_0$ を決めると、
\[
  a_0, f(a_0), f(f(a_0)),\ldots
\]
という列ができる。

この列を、関数 $f$ によって動かされる点の軌道と見る。
数学では、このように状態の更新規則を調べる分野を力学系という。
物理の運動を扱わなくても、値が次々に変化するなら力学系として考えられる。

この視点では、一般項よりも次の問いが中心になる。

\begin{itemize}
  \item 軌道はある値へ近づくか。
  \item ある範囲から飛び出すか。
  \item 周期的に繰り返すか。
  \item 初期値を少し変えると、軌道は大きく変わるか。
\end{itemize}

\subsection{不動点と周期点}

$f(L)=L$ を満たす点 $L$ を不動点という。
不動点から出発した軌道は、
\[
  L\longmapsto L\longmapsto L\longmapsto\cdots
\]
と動かない。

一方、$f^{\,p}(L)=L$ を満たし、$p$ 回より少ない反復では元に戻らない点を
周期 $p$ の周期点という。例えば、
\[
  f(x)=2-x
\]
では $1$ が不動点であり、$x\neq1$ ならば
\[
  x\longmapsto 2-x\longmapsto x
\]
となるので、周期 $2$ の軌道になる。

不動点を探すだけでは、数列の全体の動きは分からない。
不動点へ近づく軌道もあれば、そこを飛び越えて周期軌道になるものもある。

\subsection{一次関数による反復を分類する}

最も基本的な反復として、
\[
  a_{n+1}=ra_n+b
\]
を考える。$r\neq1$ のときの不動点は、
\[
  L=rL+b
\]
より、
\[
  L=\frac{b}{1-r}
\]
である。そこで $u_n=a_n-L$ と置くと、
\[
  u_{n+1}=a_{n+1}-L
  =ra_n+b-L
  =r(a_n-L)=ru_n
\]
となる。つまり、不動点からのずれは等比数列になる。
\[
  a_n-L=r^n(a_0-L)
\]
である。

この一行から、動きが分類できる。

\begin{itemize}
  \item $|r|<1$ なら $a_n\to L$。
  \item $r>1$ なら、$a_0\neq L$ の場合は不動点から遠ざかる。
  \item $r<-1$ なら、符号を交互に変えながら大きくなる。
  \item $r=-1$ なら、$a_0\neq L$ の場合は周期 $2$ になる。
  \item $r=1$ なら、$a_{n+1}=a_n+b$ であり、等差数列になる。
\end{itemize}

漸化式を解く方法と、反復の動きを分類する方法が一致した。

\subsection{分数一次変換と不動点}

\subsubsection{一次分数型は関数の反復である}

\[
  a_{n+1}=\frac{pa_n+q}{ra_n+s}
\]
を考える。ただし、各段階で分母は $0$ でないものとする。
この写像が恒等的に定数になる場合を除き、これは一次分数変換、またはメビウス変換と呼ばれる。

不動点は、
\[
  \alpha=\frac{p\alpha+q}{r\alpha+s}
\]
より、
\[
  r\alpha^2+(s-p)\alpha-q=0
\]
を満たす。分数を含む漸化式でも、不動点は二次方程式によって求められる。

\subsubsection{二つの不動点を基準にする}

異なる二つの不動点を $\alpha,\beta$ とする。次の変数変換を考える。
\[
  b_n=\frac{a_n-\alpha}{a_n-\beta}.
\]

変換
\[
  f(x)=\frac{px+q}{rx+s}
\]
について、固定点の条件を使うと、
\[
  f(x)-\alpha
  =\frac{(p-r\alpha)(x-\alpha)}{rx+s}
\]
および
\[
  f(x)-\beta
  =\frac{(p-r\beta)(x-\beta)}{rx+s}
\]
となる。したがって、
\[
  \frac{f(x)-\alpha}{f(x)-\beta}
  =\frac{p-r\alpha}{p-r\beta}
   \frac{x-\alpha}{x-\beta}.
\]

ここで
\[
  \lambda=\frac{p-r\alpha}{p-r\beta}
\]
と置けば、
\[
  b_{n+1}=\lambda b_n
\]
となる。分数を含む複雑な漸化式が、変数変換によって等比数列になった。

第2章で扱った一次分数型の解法は、二つの不動点からの比を記録する方法として理解できる。

\subsection{多項式反復と表現の変更}

\subsubsection{三角関数による翻訳}

\[
  a_{n+1}=2a_n^2-1,qquad a_1=\cos\theta
\]
を考える。倍角公式
\[
  2\cos^2t-1=\cos(2t)
\]
から、
\[
  a_n=\cos(2^{n-1}\theta)
\]
となる。

ここで行っているのは、二次式を展開することではない。
数列の値を角度の関数として表し、二次式の反復を角度の二倍へ翻訳している。

$a_1=\cosh t\geq1$ の場合には、
\[
  2\cosh^2t-1=\cosh(2t)
\]
を使って、
\[
  a_n=\cosh(2^{n-1}t)
\]
と表せる。

\subsubsection{対称式による翻訳}

三角関数を使わずに、
\[
  (u+u^{-1})^2-2=u^2+u^{-2}
\]
という恒等式を使うこともできる。
$a_1=u+u^{-1}$ で
\[
  a_{n+1}=a_n^2-2
\]
なら、
\[
  a_n=u^{2^{n-1}}+u^{-2^{n-1}}
\]
である。

置換が有効かどうかは、次の反復でも同じ表現が保たれるかによって判断できる。
表現が閉じなければ、別の方法へ戻る。

\subsection{ロジスティック写像}

次の漸化式を考える。
\[
  a_{n+1}=r a_n(1-a_n).
\]
通常、人口割合として解釈する場合は $0\leq a_0\leq1$、$0<r\leq4$ を考える。

不動点を求めると、
\[
  L=rL(1-L)
\]
より、
\[
  L=0,qquad L=1-\frac1r
\]
である。後者における微分係数は、
\[
  f'(L)=r(1-2L)=2-r
\]
となる。したがって、
\[
  |2-r|<1,qquad 1<r<3
\]
の範囲では、正の不動点が局所的に安定になる。

$r$ が $3$ を越えると、周期 $2$、周期 $4$ などの周期倍化が現れる。
この一行の漸化式は、一般項を求める問題を超えて、安定性、周期性、初期値への依存性を調べる対象になる。

\subsection*{この節のまとめ}

反復型漸化式では、次の三つの翻訳が重要である。

\begin{enumerate}
  \item 不動点からのずれへ移す。
  \item 二つの不動点からの比へ移す。
  \item 三角関数や対称式によって、非線形な反復を単純な反復へ移す。
\end{enumerate}

どの方法を選ぶかは、式の名前ではなく、置換後も同じ表現が保たれるかによって決まる。
